\section{Discussion}\label{sec:discussion}

\subsection{What is established and what is not}

Four statements are supported by the evidence, with the qualifications given. First, adaptive conformal intervals (ACI, FACI, conformal PID, DRACP) were better calibrated than a conformal calibration performed once on the first months and never updated. The direction is the same in the exploratory data, on the first hold-out panel at both warm-ups, for three base forecasters, and after removing any single country (Table~\ref{tab:loo}). The strength of the evidence depends on the warm-up size: the pre-specified rule R1 was met on the first hold-out panel only at warm-up 96, which is a sensitivity run, and at the primary warm-up 48 the confidence intervals of the contrast include zero (\ci{-5.4}{+0.5} at $h=3$ and \ci{-3.9}{+1.2} at $h=6$). The statement is therefore one of consistent direction, not of a pre-specified confirmation on that panel. On the second hold-out panel R1 was met at warm-up 48 for all four base models and not at warm-up 96, which adds a replication at the primary warm-up on new countries. Second, in a secondary, post hoc analysis, ACI with an ARIMA base model had a lower calibration error than the native ARIMA intervals at every horizon in both samples; on the first hold-out panel the interval includes zero at $h=6$ and $h=12$. Third, at the 12-month horizon pooled coverage is below nominal for every method on the design and first hold-out panels, and delayed feedback accounts for part of the shortfall, the size depending on the method and sample. Fourth, at short horizons ACI, DRACP and FACI have a median interval width below that of a rolling window in both samples, but the pooled mean width of ACI at $h=3$ is larger than that of rolling in both samples, and the Winkler ratio to rolling is below 1 only on the hold-out panel at $h=1$ and $h=3$ (0.95 and 0.99 for ACI, with an interval at $h=3$ that includes zero); no efficiency advantage is established.

Three statements are not established, and the first of them is split by panel. (i) Whether ACI improves on a rolling window depends on the panel. On the design panel and the first hold-out panel it was not shown (R2); the point estimates are negative, in the direction of an advantage of about 1 to 3~pp of CE, but they do not reach the pre-specified level and their MDE is undefined, so they are inconclusive. The re-run of the earlier 20 countries from 2010-01 shows the rule as not significant (Holm $p$ 0.22 to 1.0), again with the MDE undefined in every cell and so inconclusive under E5. On the second hold-out panel the rule is met at warm-up 48 at $h=3$ for LightGBM, ARIMA and naive B (Holm 0.008, 0.004, 0.004) and at $h=6$ for ARIMA and naive B (0.030, 0.039); for Chronos at $h=3$ the Holm $p$ is 0.036 but the MDE is undefined, so it is inconclusive under E5 and not counted as met, and Chronos and LightGBM at $h=6$ are not met (0.088, 0.159). The reason for the difference is not known, and candidate explanations were not tested. The panels are not pooled; the statement supported by the evidence taken together is that the result is panel dependent, and it is neither ``ACI does not improve on a rolling window'' nor ``ACI improves on a rolling window''. (ii) An advantage of DRACP over ACI was not detected, and equivalence within 2~pp was not stable across warm-up sizes (R3); on the second panel DRACP had a lower CE than ACI in every cell at warm-up 48 by 0.5 to 1.8~pp, with no difference at Holm $p$ below 0.05, and the verdict stays inconclusive. (iii) That delayed feedback explains the $h=12$ shortfall was not met at the pre-specified threshold (R4). None of these results is evidence that the effects are absent. In particular, the wording ``DRACP does not transfer'' or ``ACI and DRACP are equal'' is not allowed by the data. Two descriptive findings of the extension carry no inferential weight: Chronos used zero-shot as a base forecaster leaves the calibration conclusions unchanged, and on the 22 target months from 2025-01 its mean absolute error was at least as low as that of LightGBM and ARIMA. The language-model explainer was evaluated on one model with programmatic metrics (Section~\ref{sec:explainer}); that result is outside the main claims, and nothing general is claimed about the faithfulness of such a layer: the grounded result is a ceiling of what the metrics see, the corrected metrics are post hoc and were checked by one reader that was the model, and the ungrounded condition is a floor.

\subsection{Why null results: power, the ceiling effect and the equivalence margin}\label{sec:nulls}

Several reasons may contribute to the null outcomes, and the data do not allow us to apportion them.

\paragraph{Sample size and dependence.} The first hold-out panel has 7 countries and about 200 origin months at $h=3$. Errors of $h$-step forecasts overlap, adaptive methods have a long memory (about 50 months at step 0.02), and countries share common shocks. The block bootstrap respects this dependence at the cost of wide intervals; the 95 percent intervals of the primary contrasts are about 2 to 6~pp wide on the CE scale.

\paragraph{Ceiling effect.} The calibration error of the rolling baseline is already small on the first hold-out panel: 0.029 at $h=1$, 0.036 at $h=3$ and 0.032 at $h=6$, and ACI's is 0.020, 0.020 and 0.025. There is little room to improve. The power simulation reflects this: its alternative cannot reduce CE below zero, so the improvement it can create is bounded by the baseline CE (about 3 to 4~pp), and 80 percent power at the first Holm step was not reached up to 10~pp on the full hold-out sample. The MDE is therefore undefined for all four primary contrasts (R5). The same is true for the exploratory sample. On random 15-month windows the median MDE is 2.6 to 3.1~pp on the hold-out panel, but it is undefined in 63 to 68 percent of windows. A contrast without defined power cannot support a claim of absence of an effect, and the pre-specified R5 therefore classifies all four contrasts as inconclusive. We interpret this as a limit of the design (a small and already well-calibrated reference) and not as a finding about the methods.

\paragraph{The equivalence margin is about the size of ACI's own calibration error.} The margin of 2~pp is about the size of the calibration error of ACI itself (0.013 and 0.014 at $h=1$ and $h=3$ in the exploratory data, and 0.020 and 0.020 on the hold-out panel). A test of whether DRACP differs from ACI by less than 2~pp therefore asks whether a method with an error of about 2~pp differs from another by less than that amount, and the verdict ``equivalent'' is weak by construction. For the same reason the contrasts against ACI cannot be powered: this is a ceiling effect, and the null outcomes for ACI against rolling and for DRACP against ACI were likely before the run. The informative contrasts are those against weak references, split (CE 0.050 to 0.056 on the hold-out panel at $h=1$ to 6) and the native intervals, where there is room for a difference. The margin was fixed in advance and is not justified beyond that.

\paragraph{The ACI-against-rolling result depends on the panel.} The second hold-out panel is an honest counter-example to the null outcomes above, and it must be reported next to them. At warm-up 48 the pre-specified rule R2 is met there in part of the cells (Table~\ref{tab:r2panels}), where it was not shown on the design panel and the first hold-out panel, and it is not significant on the earlier 20 countries re-run from 2010-01, where the MDE is undefined in every cell, so those results are inconclusive and not evidence of no effect. The data do not say why. Panel differences are candidate explanations: a shorter series (198 months from 2010-01), lower-income and more volatile economies, a warm-up that takes a larger share of 198 months, and a larger calibration error of the rolling window, which leaves more room (0.066 to 0.090 at $h=3$ against 0.036 on the first hold-out panel). The last point is consistent with the ceiling argument above, because on this panel the minimum detectable difference for ACI minus rolling is defined (1.2 to 6.6~pp), but consistency is not a test, and none of these candidates was tested. Warm-up 96 on the new panel cannot arbitrate: it has 41 and 35 pooled months, an undefined minimum detectable difference and point estimates of both signs. Two reading errors are avoided. The panels are not averaged into a single estimate, because they differ in the number of countries, the length of the series, the start date and the base models, and because the pre-specified rules apply to each panel by itself. The favourable panel is not selected either: the evidence is one panel on which the rule is met in part of the cells and two on which it is not, and a correction across the four base models, which share their countries and months, was not applied. A reader who wants a recommendation for a new set of countries cannot take it from this study.

\paragraph{A true small or zero difference.} The data are also compatible with ACI and a rolling window being close to each other in this setting, and with DRACP and ACI being close, but this cannot be separated from low power with the present sample.

\paragraph{Specific to DRACP.} DRACP trains its regime and scale models on half of the warm-up (24 rows at warm-up 48), which is a small sample for these models, and its results depend on the warm-up: equivalence holds at 48 and not at 96. In the exploratory data at warm-up 96, DRACP was wider than ACI: its Winkler score was 33 percent higher at $h=3$. This does not hold on the first hold-out panel: at warm-up 96 the Winkler score of DRACP was 5.8 percent lower than that of ACI at $h=3$ (95\% interval \ci{-14.2}{+3.3}, unadjusted $p=0.18$), with a mean width of 6.11 against 6.31, and the mean width of DRACP is larger than ACI's only in the exploratory data (4.97 against 4.28 at $h=3$). The authors of DRACP state that the advantage of their method is specific to one-step-ahead forecasting and depends on the point forecaster; the present design uses multi-month cumulative targets, a different point forecaster and a different panel, and it does not test their one-step claim. Our reproduction of the one-step setting was partial and qualitative: ACI 0.874 and DRACP 0.871 in coverage over 5 countries with one seed, against 0.884 and 0.890 reported by the authors over 48 series, so the claim of being closest to nominal could not be distinguished from ACI there either. This is a limitation of the transfer comparison and not a finding about the method.

\subsection{Calibrated is not the same as useful}

A prediction interval can reach its nominal coverage with very wide intervals, and calibration error alone does not say whether the interval is informative. This is why width and the interval score are reported next to CE. The widths show the cost: in the exploratory data the pooled mean width of ACI intervals is 4.28 percentage points of cumulative growth at $h=3$ and 15.5 at $h=12$, and at long horizons the better calibration of the adaptive methods is paid for in width or interval score. At $h=12$ an interval with honest 90 percent coverage may be too wide for practical policy or business use. The study does not establish a width at which an interval becomes useful and does not compare widths with a decision-relevant threshold.

\subsection{Implications for practice}

A practitioner who produces inflation intervals for a commodity-dependent economy should not calibrate once and leave the calibration unchanged: one-shot calibration had a higher calibration error than ACI in every CE contrast reported with LightGBM, ARIMA and naive B (both samples, both warm-ups), and for Chronos in all but one cell (the second hold-out panel at warm-up 96 and $h=6$, where the difference is $+0.8$~pp and not distinguishable from zero), although the pre-specified rule R1 was met on the first hold-out panel only at the sensitivity warm-up and the confidence intervals at the primary warm-up include zero. Its coverage also varied the most across countries (the standard deviation of coverage across units at $h=12$ is 0.152 for split against 0.058 for ACI on the first hold-out panel). The native intervals of an ARIMA model had a higher calibration error than ACI on the same model in both samples, in a post hoc comparison. Updating the calibration, with a rolling window or with ACI, had a lower calibration error than calibrating once in every cell reported (a descriptive comparison for the rolling window, for which no contrast against split was tested) and is cheap to implement. Between a rolling window and ACI this study cannot recommend one on statistical grounds, because the comparison depends on the panel (not shown on two panels, met at warm-up 48 in part of the cells on a third), and it cannot recommend one on efficiency grounds either: the median width of ACI is smaller at short horizons, but the pooled mean width and the Winkler score do not show a consistent gain. At the 12-month horizon pooled coverage is below nominal for every method (0.83 to 0.88 on the first hold-out panel), and users should treat the nominal level there as an upper bound. A part of the shortfall comes from the delay with which the outcome arrives and cannot be removed in real use.

\subsection{The choice of ACI in the platform}

The platform uses ACI. ACI was chosen as an adaptive method that is better calibrated than a one-shot calibration and than native intervals, and it is simple to implement and explain. The benchmark did not show it to be better than a rolling window on the design panel and the first hold-out panel, and on the second hold-out panel it was better at warm-up 48 in part of the cells; the result depends on the panel. No efficiency advantage is claimed. The platform module is an engineering application with a simpler point forecast and is not a scientific contribution of the paper. It displays measured backtest coverage next to the nominal level so that a user can see the actual calibration.

\subsection{Scope of the hold-out and of the term ``commodity-dependent''}

The hold-out countries are independent of the design sample by country only. They share calendar months with it, including the global shocks of 2020 to 2022, so common shocks affect both samples. The methods, the contrasts and the horizons were chosen on the design countries. The agreement between the samples is therefore weaker evidence of generality than if the hold-out countries had been observed in different periods. The second hold-out panel (2010 to 2026) shares calendar months with both earlier samples in the same way, and its rules E1--E5 were written after all earlier results were known. Its agreement with the earlier samples on R1 and its disagreement on R2 are both results about new countries and not about new periods. No export-share criterion defines commodity dependence in any panel, so the phrase in the title describes the intended setting and not a tested property.
